% GENERATED FILE. Edit canonical lessons, then run npm run generate:books.
% canonical-source-hash: fnv1a64:9d582c8e6ee0f76a
% canonical-lessons: KA-C244-sarisu, KA-C244-kuggu, KA-C244-nekku, KA-C244-hiru, KA-C244-alavadisu

\chapter{To move aside, To shrink, To lick, To suck, To install}
\label{ch:ka-to-move-aside-to-shrink-to-lick-to-suck-to-install}
\hlchaptermodality{244}

\begin{chapteropening}
\textbf{By the end of this chapter:} \emph{I can say to move aside, to shrink, to lick, to suck and to install.}

Complete the last lesson of chapter 244: I can say to move aside, to shrink, to lick, to suck and to install.
\end{chapteropening}

\section[\texorpdfstring{\kn{ಸರಿಸು} (sarisu)}{ಸರಿಸು (sarisu)}]{\kn{ಸರಿಸು} (sarisu) — to move aside, to shift}

\label{lesson:KA-C244-sarisu}



\begin{quote}
Before the new one: say the Kannada for to manage, then the Kannada for to shake.
\end{quote}

\subsection*{You'll want to know: \kn{ಸರಿಸು}}
\textbf{\kn{ಸರಿಸು}} — \emph{sarisu} — \textquotedblleft{}to move aside, to shift\textquotedblright{}.

\begin{grammarlens}[title={What you've built}]
One more word for this chapter.
\end{grammarlens}

\subsection*{Your turn}
\begin{itemize}
\raggedright
  \item \emph{Say it:} \emph{sarisu}
  \item \emph{Say it:} \emph{sarisu}, once more
  \item \emph{Say it:} say \emph{sarisu}
  \item \emph{Recall:} say the Kannada for to manage, then the Kannada for to shake, then say \emph{sarisu} again
\end{itemize}

\begin{tcolorbox}[breakable,colback=teal!4,colframe=teal!35!black,title={Before you move on}]
What does \kn{ಸರಿಸು} mean? (\textquotedblleft{}To move aside, to shift\textquotedblright{}.) Say it once more.
\end{tcolorbox}
\section[\texorpdfstring{\kn{ಕುಗ್ಗು} (kuggu)}{ಕುಗ್ಗು (kuggu)}]{\kn{ಕುಗ್ಗು} (kuggu) — to shrink}

\label{lesson:KA-C244-kuggu}



\begin{quote}
Before the new one: say the Kannada for to shake, then the Kannada for to move aside.
\end{quote}

\subsection*{You'll want to know: \kn{ಕುಗ್ಗು}}
\textbf{\kn{ಕುಗ್ಗು}} — \emph{kuggu} — \textquotedblleft{}to shrink\textquotedblright{}.

\begin{grammarlens}[title={What you've built}]
One more word for this chapter.
\end{grammarlens}

\subsection*{Your turn}
\begin{itemize}
\raggedright
  \item \emph{Say it:} \emph{kuggu}
  \item \emph{Say it:} \emph{kuggu}, once more
  \item \emph{Say it:} say \emph{kuggu}
  \item \emph{Recall:} say the Kannada for to shake, then the Kannada for to move aside, then say \emph{kuggu} again
\end{itemize}

\begin{tcolorbox}[breakable,colback=teal!4,colframe=teal!35!black,title={Before you move on}]
What does \kn{ಕುಗ್ಗು} mean? (\textquotedblleft{}To shrink\textquotedblright{}.) Say it once more.
\end{tcolorbox}
\section[\texorpdfstring{\kn{ನೆಕ್ಕು} (nekku)}{ನೆಕ್ಕು (nekku)}]{\kn{ನೆಕ್ಕು} (nekku) — to lick}

\label{lesson:KA-C244-nekku}



\begin{quote}
Before the new one: say the Kannada for to move aside, then the Kannada for to shrink.
\end{quote}

\subsection*{You'll want to know: \kn{ನೆಕ್ಕು}}
\textbf{\kn{ನೆಕ್ಕು}} — \emph{nekku} — \textquotedblleft{}to lick\textquotedblright{}.

\begin{grammarlens}[title={What you've built}]
One more word for this chapter.
\end{grammarlens}

\subsection*{Your turn}
\begin{itemize}
\raggedright
  \item \emph{Say it:} \emph{nekku}
  \item \emph{Say it:} \emph{nekku}, once more
  \item \emph{Say it:} say \emph{nekku}
  \item \emph{Recall:} say the Kannada for to move aside, then the Kannada for to shrink, then say \emph{nekku} again
\end{itemize}

\begin{tcolorbox}[breakable,colback=teal!4,colframe=teal!35!black,title={Before you move on}]
What does \kn{ನೆಕ್ಕು} mean? (\textquotedblleft{}To lick\textquotedblright{}.) Say it once more.
\end{tcolorbox}
\section[\texorpdfstring{\kn{ಹೀರು} (hīru)}{ಹೀರು (hīru)}]{\kn{ಹೀರು} (hīru) — to suck, to sip}

\label{lesson:KA-C244-hiru}



\begin{quote}
Before the new one: say the Kannada for to shrink, then the Kannada for to lick.
\end{quote}

\subsection*{You'll want to know: \kn{ಹೀರು}}
\textbf{\kn{ಹೀರು}} — \emph{hīru} — \textquotedblleft{}to suck, to sip\textquotedblright{}.

\begin{grammarlens}[title={What you've built}]
One more word for this chapter.
\end{grammarlens}

\subsection*{Your turn}
\begin{itemize}
\raggedright
  \item \emph{Say it:} \emph{hīru}
  \item \emph{Say it:} \emph{hīru}, once more
  \item \emph{Say it:} say \emph{hīru}
  \item \emph{Recall:} say the Kannada for to shrink, then the Kannada for to lick, then say \emph{hīru} again
\end{itemize}

\begin{tcolorbox}[breakable,colback=teal!4,colframe=teal!35!black,title={Before you move on}]
What does \kn{ಹೀರು} mean? (\textquotedblleft{}To suck, to sip\textquotedblright{}.) Say it once more.
\end{tcolorbox}
\section[\texorpdfstring{\kn{ಅಳವಡಿಸು} (aḷavaḍisu)}{ಅಳವಡಿಸು (aḷavaḍisu)}]{\kn{ಅಳವಡಿಸು} (aḷavaḍisu) — to install, to fit}

\label{lesson:KA-C244-alavadisu}



\begin{quote}
Before the new one: say the Kannada for to lick, then the Kannada for to suck.
\end{quote}

\subsection*{You'll want to know: \kn{ಅಳವಡಿಸು}}
\textbf{\kn{ಅಳವಡಿಸು}} — \emph{aḷavaḍisu} — \textquotedblleft{}to install, to fit\textquotedblright{}.

\begin{grammarlens}[title={What you've built}]
One more word for this chapter.
\end{grammarlens}

\subsection*{Your turn}
\begin{itemize}
\raggedright
  \item \emph{Say it:} \emph{aḷavaḍisu}
  \item \emph{Say it:} \emph{aḷavaḍisu}, once more
  \item \emph{Say it:} say \emph{aḷavaḍisu}
  \item \emph{Recall:} say the Kannada for to lick, then the Kannada for to suck, then say \emph{aḷavaḍisu} again
\end{itemize}

\begin{tcolorbox}[breakable,colback=teal!4,colframe=teal!35!black,title={Before you move on}]
What does \kn{ಅಳವಡಿಸು} mean? (\textquotedblleft{}To install, to fit\textquotedblright{}.) Say it once more.
\end{tcolorbox}
